% !TEX root = ../../main.tex
% C4.4 - Use case diagram (translated). Figure H1 is the English version (figures/H1-use-case.png, generated from tools-en).
\section{Use Case Diagram}
\label{sec:4.4}

The functional requirements of Section~\ref{sec:4.2} are organised into 15 use cases. Figure~\ref{fig:usecase} shows all the use cases of the system with their relations and actors, and Table~\ref{tab:usecases} lists each use case with its actor and purpose.

\begin{figure}[H]
\centering
\includegraphics[width=\textwidth]{H1-use-case.png}
\caption{Use case diagram of the system}
\label{fig:usecase}
\end{figure}

The diagram is organised as follows:
\begin{itemize}
    \item \textbf{An abstract \emph{User} actor.} The three concrete actors share login (UC-01) and logout (UC-15), so the diagram links these two use cases to an abstract \emph{User} actor, which the three concrete actors generalise to.
    \item \textbf{Division by role.} The Administrator is linked to UC-02 to UC-08, the Operator to UC-09 to UC-11 and the Viewer to UC-12 to UC-14; no use case is shared among these actors.
    \item \textbf{Extend relations.} UC-10 extends UC-09 and UC-11 extends UC-10: after searching, the Operator can view the details of a result and then save it into a case. Likewise, UC-13 extends UC-12 and UC-14 extends UC-13. The extending use cases are still linked directly to their actor, because the actor can also start them independently, for example opening the list of cases without searching first.
    \item \textbf{Login is a precondition} of every other use case, not a step inside them, so the diagram draws no «include» relation to UC-01.
\end{itemize}

The automatic system functions (Section~\ref{sec:4.2.5}) have no use case of their own because no actor triggers them; they take place when the Administrator enables AI processing in UC-04.

\begin{table}[H]
\centering
\caption{List of use cases}
\label{tab:usecases}
\begin{tabularx}{\textwidth}{|>{\raggedright\arraybackslash}p{1.4cm}|>{\raggedright\arraybackslash}p{3.8cm}|>{\raggedright\arraybackslash}p{2.4cm}|L|}
\hline
\textbf{Code} & \textbf{Use case} & \textbf{Actor} & \textbf{Purpose} \\
\hline
UC-01 & Log in & All three actors & Authenticate and access the functions of the role \\
\hline
UC-02 & Manage user accounts & Administrator & Create, update, lock and delete accounts; assign roles and areas \\
\hline
UC-03 & Manage cameras & Administrator & Add and update cameras, check the RTSP connection, retire and restore them \\
\hline
UC-04 & Manage AI processing of cameras & Administrator & Enable or disable AI processing per camera; process fallback video files \\
\hline
UC-05 & Configure AI models & Administrator & Choose the Detector and Tracker used by the whole system \\
\hline
UC-06 & Monitor system status & Administrator & Follow the status of cameras, RTSP streams and AI processing \\
\hline
UC-07 & Check AI operation & Administrator & Check the camera processing group and the search group \\
\hline
UC-08 & View the audit log & Administrator & Consult the log of actions and technical events \\
\hline
UC-09 & Search for people with AI & Operator & Search by image, description or attributes within the assigned area \\
\hline
UC-10 & View and assess search results & Operator & View the details of each result, compare and assess it \\
\hline
UC-11 & Manage case files & Operator & Create cases, save results, update, close and reopen cases \\
\hline
UC-12 & View the overview & Viewer & See the system-wide figures and recent cases \\
\hline
UC-13 & View case files & Viewer & View every case in read-only mode \\
\hline
UC-14 & View saved search results & Viewer & View the details of a result saved in a case \\
\hline
UC-15 & Log out & All three actors & End the session \\
\hline
\end{tabularx}
\end{table}
